\chapter{Conclusion}
\label{chap:conclusion}
Aegis set out to answer ``where did I leave it?'' for a person with memory
impairment from first-person video. At the end of FYP-1 the whole path exists: a
phone streams video, a laptop keeps sharp, new frames and uploads segments,
Modal reconstructs them in 3D and detects objects, the memory records where each
object is relative to furniture and rooms in a bounded, persistent per-home
store, and two apps let the patient and caregiver ask by tapping, typing or
speaking, with a photo as evidence. 451 tests pass with 79\% coverage.

Three results stand out. SAM~3 made small personal objects detectable (F1 0.721
against 0.097). State runs made memory bounded (5.7--7.2$\times$ smaller with no
loss of last-known states). And the store kept every tracked object across days
on EgoLife.

The report is equally clear about what is not done. Temporal answers,
calibrated confidence, authentication and encryption are not built; accuracy,
latency and usability are not measured; and the prototype is not safe for
patient data. FYP-2 (Chapter~\ref{chap:fyp2}) closes these gaps and measures
the results before any claim is made.
